\title{The Era of 1-bit LLMs: All Large Language Models are in 1.58 Bits}

\begin{abstract}
Recent advancements, exemplified by BitNet~\cite{bitnet}, are ushering in a new generation of 1-bit Large Language Models (LLMs). In this study, we present a 1-bit LLM variant, referred to as \textbf{\text{BitNet b1.58}}, where each parameter (or weight) of the LLM is ternary \{-1, 0, 1\}. This model achieves parity with full-precision (i.e., FP16 or BF16) Transformer LLMs of equivalent model size and training data in terms of both perplexity and downstream task performance, while offering substantial improvements in latency, memory usage, throughput, and energy efficiency. More importantly, the 1.58-bit LLM establishes a novel scaling law and methodology for developing future LLM generations that balance high performance and cost-effectiveness. Additionally, it introduces a new computational paradigm and lays the groundwork for creating specialized hardware tailored to 1-bit LLMs.
\end{abstract}

\section{The Era of 1-bit LLMs}

In recent years, the AI community has witnessed rapid expansion in both the size and capabilities of Large Language Models (LLMs). These models have exhibited outstanding results across a broad spectrum of natural language processing tasks; however, their growing size has introduced challenges for deployment and sparked concerns about environmental and economic impacts due to substantial energy consumption.

One strategy to tackle these issues is post-training quantization, which produces low-bit models for inference~\cite{smoothquant, gptq, quip, quip_sharp}. This method decreases the precision of weights and activations, thereby significantly lowering memory and computational demands of LLMs. The prevailing trend has shifted from 16-bit representations to even lower bit widths, such as 4-bit models~\cite{gptq, awq}. Nonetheless, post-training quantization remains a suboptimal solution, despite its widespread adoption in industry-scale LLMs.

Recent studies on 1-bit model architectures, like BitNet~\cite{bitnet}, offer a promising approach to decreasing the cost of LLMs while preserving their effectiveness. Standard LLMs utilize 16-bit floating-point values (i.e., FP16 or BF16), with most of their operations dominated by matrix multiplication. Consequently, the primary computational expense arises from floating-point addition and multiplication. In contrast, BitNet’s matrix multiplication relies solely on integer addition, leading to substantial energy savings for LLMs. Since power consumption is often the limiting factor in chip performance, these energy savings can translate into faster computations.

Beyond computation, transferring model parameters from DRAM to the memory of an on-chip accelerator (such as SRAM) can incur significant costs during inference. Although attempts have been made to expand SRAM to boost throughput, this approach incurs far greater expenses than using DRAM. Compared to full-precision models, 1-bit LLMs have a considerably smaller memory footprint in terms of both capacity and bandwidth. This reduction greatly decreases the cost and time required to load weights from DRAM, enabling quicker and more efficient inference.

In this study, we propose a notable 1-bit LLM variant named \textbf{\text{BitNet b1.58}}, where every parameter is ternary, adopting values from the set \{-1, 0, 1\}. By introducing an additional zero value to the original 1-bit BitNet, we achieve 1.58 bits in the binary system. \text{BitNet b1.58} preserves all the advantages of the original 1-bit BitNet, including its novel computation approach that nearly eliminates multiplication operations during matrix multiplication and can be highly optimized. Moreover, it maintains the same energy consumption as the original 1-bit BitNet and significantly surpasses FP16 LLM baselines in memory efficiency, throughput, and latency. Additionally, \text{BitNet b1.58} provides two further benefits. First, its modeling power is enhanced by explicitly supporting feature filtering through the inclusion of zero in the model weights, which notably boosts the performance of 1-bit LLMs. Second, our results demonstrate that \text{BitNet b1.58} can achieve parity with full-precision (i.e., FP16) baselines in both perplexity and end-task outcomes, starting from a 3B parameter size, under identical configurations (e.g., model size, training tokens, etc.).

\section{BitNet b1.58}

\text{BitNet b1.58} is built upon the BitNet framework, which modifies the Transformer by substituting \emph{nn.Linear} layers with \emph{BitLinear}. It is trained from the beginning, utilizing weights quantized to 1.58 bits and activations quantized to 8 bits. Compared to the original BitNet, it incorporates several changes which we outline below.

\paragraph{Quantization Function.} To restrict the weights to values in \{-1, 0, +1\}, we employ an \emph{absmean} quantization method. This approach initially normalizes the weight matrix by dividing by its average absolute value, then rounds each entry to the closest integer in \{-1, 0, +1\}:

\begin{equation}
    \widetilde{W} = \mathrm{RoundClip}\left(\frac{W}{\gamma + \epsilon}, -1, 1\right),
\end{equation}

\begin{equation}
    \mathrm{RoundClip}(x, a, b) = \max\bigl(a, \min(b, \mathrm{round}(x))\bigr),
\end{equation}

\begin{equation}
    \gamma = \frac{1}{nm}\sum_{ij} |W_{ij}|.
\end{equation}

For activation quantization, we follow the same procedure as in BitNet, except that we do not scale the activations prior to the non-linearities to the range $[0, Q_b]$. Instead, each token’s activations are scaled within the range $[-Q_b, Q_b]$ to avoid zero-point quantization. This adjustment simplifies the implementation and system-level optimization while having negligible impact on performance in our tests.

\paragraph{LLaMA-alike Components.}

LLaMA~\cite{llama, llama2} has become the standard backbone architecture for open-source large language models. To align with the open-source ecosystem, our \text{BitNet b1.58} design incorporates LLaMA-style elements. In particular, it utilizes RMSNorm~\cite{rmsnorm}, SwiGLU~\cite{swiglu}, rotary embeddings~\cite{rope}, and removes all bias terms. This design choice enables \text{BitNet b1.58} to be easily integrated into widely used open-source frameworks (such as Huggingface, vLLM~\cite{vllm}, and llama.cpp\footnote{\url{https://github.com/ggerganov/llama.cpp}}) with minimal modifications.

\section{Results}

We conducted a comparison between \text{BitNet b1.58} and our reimplemented FP16 \text{LLaMA LLM} across multiple model sizes. To maintain fairness, both models were pre-trained on the RedPajama dataset~\cite{redpajama} with 100 billion tokens. We assessed their zero-shot performance on a variety of language tasks, such as ARC-Easy~\cite{arc}, ARC-Challenge~\cite{arc}, Hellaswag~\cite{hellaswag}, Winogrande~\cite{winoGrande}, PIQA~\cite{piqa}, OpenbookQA~\cite{openbookqa}, and BoolQ~\cite{boolq}. Additionally, we measured validation perplexity using the WikiText2~\cite{wikitext2} and C4~\cite{c4} datasets.

We also evaluated the runtime GPU memory consumption and inference latency for both \text{LLaMA LLM} and \text{BitNet b1.58}. These measurements were performed using the FasterTransformer\footnote{\url{https://github.com/NVIDIA/FasterTransformer}} framework, which is highly optimized for LLM inference latency on GPU platforms. The 2-bit kernel from Ladder~\cite{ladder} was incorporated for \text{BitNet b1.58}. The reported metric was the time taken per output token, as this represents the principal computational cost during inference.

Table~\ref{tab:ppl} provides a summary of the perplexity scores and computational costs for both \text{BitNet b1.58} and \text{LLaMA LLM}. The results indicate that \text{BitNet b1.58} attains comparable perplexity to the full precision \text{LLaMA LLM} starting at the 3B model scale, while achieving 2.71 times faster inference and requiring 3.55 times less GPU memory. Notably, the 3.9B \text{BitNet b1.58} model operates 2.4 times faster and consumes 3.32 times less memory, yet significantly outperforms the \text{LLaMA LLM} 3B model.

Detailed zero-shot accuracy results on downstream tasks are presented in Table~\ref{tab:zero-shot}. The evaluation was carried out following the methodology of the \emph{lm-evaluation-harness}\footnote{\url{https://github.com/EleutherAI/lm-evaluation-harness}}. Findings show the performance gap between \text{BitNet b1.58} and \text{LLaMA LLM} diminishes as model size increases. Crucially, \text{BitNet b1.58} matches the full precision baseline’s performance starting at the 3B scale. Consistent with the perplexity observations, the downstream task results demonstrate that \text{BitNet b1.58} 3.9B surpasses \text{LLaMA LLM} 3B while requiring less memory and latency, establishing \text{BitNet b1.58} as a Pareto improvement relative to current state-of-the-art LLMs.

\paragraph{Memory and Latency}

We expanded the model size to 7B, 13B, and 70B and assessed the associated costs. Figure~\ref{fig:latency-memory-energy} displays the patterns in latency and memory usage, indicating that the speed-up improves as the model size increases. Notably, \text{BitNet b1.58} 70B achieves a speed that is 4.1 times greater than the \text{LLaMA LLM} baseline. This improvement arises because the time required for \emph{nn.Linear} operations increases with the model size. Memory usage follows a similar pattern, since the embedding remains in full precision and its share of memory decreases in larger models. Both latency and memory metrics were obtained using a 2-bit kernel, suggesting further optimizations are possible to reduce costs even more.

\paragraph{Energy}

We also evaluate the energy consumption of arithmetic operations for both \text{BitNet b1.58} and \text{LLaMA LLM}. Our main focus is on matrix multiplication calculations, as they dominate the computational expense in LLMs. Figure~\ref{fig:energy} breaks down the energy cost composition. Most of the operations in \text{BitNet b1.58} involve INT8 addition, whereas \text{LLaMA LLM} includes both FP16 addition and FP16 multiplication. Based on the energy model from~\cite{energycost, pokebnn}, \text{BitNet b1.58} achieves a 71.4-fold reduction in arithmetic operation energy for matrix multiplications on 7nm technology. We also present end-to-end energy costs for models processing 512 tokens. Our findings indicate that as model size grows, \text{BitNet b1.58} becomes progressively more energy-efficient relative to the FP16 \text{LLaMA LLM} baseline. This trend occurs because the proportion of \emph{nn.Linear} operations increases with model size, while the energy cost from other components diminishes in larger models.

\paragraph{Throughput}

We evaluate the throughput of \text{BitNet b1.58} and \text{LLaMA LLM} with 70B parameters using two 80GB A100 GPUs, employing pipeline parallelism~\cite{gpipe} to enable running \text{LLaMA LLM} 70B on these devices. The batch size was increased until reaching the GPU memory capacity, with a fixed sequence length of 512. As shown in Table~\ref{tab:throughput}, \text{BitNet b1.58} 70B can handle a batch size up to 11 times larger than \text{LLaMA LLM}, resulting in a throughput improvement of 8.9 times.

\textbf{\text{BitNet b1.58} introduces a novel scaling law regarding model performance and inference cost}. For comparison, based on the data presented in Figures~\ref{fig:latency-memory-energy} and \ref{fig:energy}, the following equivalences between model sizes in 1.58-bit and 16-bit formats can be established:

\begin{itemize}
    \item 13B BitNet b1.58 surpasses 3B FP16 LLM in efficiency across latency, memory consumption, and energy use.
    \item 30B BitNet b1.58 outperforms 7B FP16 LLM in terms of latency, memory usage, and energy efficiency.
    \item 70B BitNet b1.58 exceeds 13B FP16 LLM in latency, memory, and energy efficiency metrics.
\end{itemize}

\paragraph{Training with 2T Tokens}

The quantity of training tokens is a key parameter for LLM performance. To assess the token scalability of \text{BitNet b1.58}, we trained a model with 2T tokens following the data methodology used for StableLM-3B~\cite{StableLM-3B-4E1T}, currently the leading open-source 3B model. Both models were tested on a benchmark comprising Winogrande~\cite{winoGrande}, PIQA~\cite{piqa}, SciQ~\cite{sciq}, LAMBADA~\cite{lambada}, and ARC-easy~\cite{arc}. The zero-shot accuracy results are reported in Table~\ref{tab:2t}. For tasks evaluated with accuracy and normalized accuracy, the average was taken. The StableLM 3B results at 2T tokens were obtained from its official technical report. Our results demonstrate that \text{BitNet b1.58} outperforms StableLM on all tasks, suggesting that 1.58-bit LLMs maintain strong generalization capabilities.

\section{Discussion and Future Work}

\textbf{1-bit Mixture-of-Experts (MoE) LLMs}

Mixture-of-Experts (MoE) has been shown to be an efficient approach for large language models (LLMs). Although it considerably lowers the computational FLOPs, its deployment and usability are constrained by substantial memory requirements and inter-chip communication overheads. These issues can be mitigated by employing 1.58-bit LLMs. Specifically, the smaller memory footprint decreases the number of devices needed to host MoE models. Additionally, it greatly diminishes the cost of transferring activations across networks. Ideally, if the entire model fits on a single chip, these overheads would be eliminated altogether.

\textbf{Native Support of Long Sequence in LLMs}

Handling long sequences has become an essential capability in the era of LLMs. A significant obstacle during long sequence inference is the memory overhead caused by KV caches. \text{BitNet b1.58} marks a notable advancement towards intrinsic support for long sequences by halving activation precision from 16 bits to 8 bits, effectively allowing the context length to double under the same resource constraints. Furthermore, these activations can be losslessly compressed down to 4 bits or even fewer for 1.58-bit LLMs, which remains an avenue for future exploration.

\textbf{LLMs on Edge and Mobile}

The deployment of 1.58-bit LLMs holds substantial promise for enhancing the performance of language models on edge and mobile platforms. These platforms typically face restrictions due to limited memory and processing power, which can hamper the scale and capabilities of LLMs. The lower memory and energy demands of 1.58-bit LLMs make their deployment on such devices feasible, unlocking a variety of applications that were previously unattainable. This advancement can significantly boost the functionalities of edge and mobile devices, paving the way for innovative LLM applications. Moreover, 1.58-bit LLMs exhibit better compatibility with CPU architectures, which are prevalent in edge and mobile hardware. Consequently, \text{BitNet b1.58} can be run efficiently on these platforms, further enhancing their overall performance and feature set.

\textbf{New Hardware for 1-bit LLMs}

Recent developments such as Groq\footnote{\url{https://groq.com/}} have showcased encouraging results and great potential in designing specialized hardware (e.g., LPUs) tailored for LLMs. Building on this progress, we advocate for the design of new hardware and systems specifically optimized for 1-bit LLMs, leveraging the novel computational paradigm introduced by BitNet~\cite{bitnet}.

\end{document}